\chapter*{Tài liệu tham khảo}

\addcontentsline{toc}{chapter}{Tài liệu tham khảo}

% Original numbering and published bibliographic titles are retained.

\begin{description}[style=nextline,leftmargin=12mm,labelwidth=10mm,font=\normalfont]

\item[{[1]}] jUCMNav User Requirements Notation Tool. \url{http://jucmnav.softwareengineering.ca/jucmnav/}.

\item[{[2]}] Lego MINDSTORMS NXT Hardware Developer Kit. Technical report, Lego, 2006.

\item[{[3]}] Complexity-Reducing Design Patterns for Cyber-Physical Systems, Steven P. Miller (presentation at the AADL standards committee meeting). \url{https://wiki.sei.cmu.edu/aadl/images/8/85/Rockwell-META-CPS--Jan-2011.pdf}, 2010.

\item[{[4]}] Competition on software verification, 2014. \url{http://sv-comp.sosy-lab.org/2014/}.

\item[{[5]}] Airbus 380 Cables Bug. \url{http://calleam.com/WTPF/?p=4700}.

\item[{[6]}] AADL Website. \url{http://www.aadl.info/aadl/}.

\item[{[7]}] AADL Inspector Homepage. \url{https://www.ellidiss.com/products/aadl-inspector/}.

\item[{[8]}] E. Allen, R. Cartwright, and B. Stoler. DrJava: A Lightweight Pedagogic Environment for Java. ACM, 2002.

\item[{[9]}] M. Andersson. Object-Oriented Modeling and Simulation of Hybrid Systems. PhD thesis, 1994.

\item[{[10]}] Ariane 5 Bug. \url{https://blog.bugsnag.com/bug-day-ariane-5-disaster/}.

\item[{[11]}] A. Asghar, S. Tariq, M. Torabzadeh-Tari, P. Fritzson, A. Pop, M. Sjölund, P. Vasaiely, and W. Schamai. An open source modelica graphic editor integrated with electronic notebooks and interactive simulation. In Proceedings of the 8th International Modelica Conference - Modelica’2011, Dresden, Germany, 2011.

\item[{[12]}] U. Assmann. Invasive Software Composition. Spinger Verlag, 2002.

\item[{[13]}] M. Association. Modelica - A Unified Object-Oriented Language for System Modelling and Simulation. 2018.

\item[{[14]}] K. J. Astrom, H. Elmqvist, D. Ab, and S. E. Mattsson. Evolution of continuous-time modeling and simulation, 1998.

\item[{[15]}] K. J. Åström, H. Elmqvist, and S. E. Mattsson. Evolution of Continuous-Time Modeling and Simulation. In ESM, pages 9–18, 1998.

\item[{[16]}] K. J. Aström and R. M. Murray. Feedback systems: an introduction for scientists and engineers. Princeton university press, 2010.

\item[{[17]}] R. B. Atitallah, E. Senn, D. Chillet, M. Lanoe, and D. Blouin. An Efficient Framework for Power-Aware Design of Heterogeneous MPSoC. IEEE Transactions on Industrial Informatics, 9(1):487–501, Feb 2013.

\item[{[18]}] B. Bachmann. Proceedings of the 6th International Modelica Conference. Bielefeld University, Bielefeld, Germany, 2008.

\item[{[19]}] M. Barbero, M. D. Del Fabro, and J. Bézivin. Traceability and provenance issues in global model management. In ECMDA-TW: Proceedings of the 3rd Workshop on Traceability, 2007.

\item[{[20]}] J. Barnat, L. Brim, and P. Ročkai. Scalable multi-core ltl model-checking. In Model Checking Software, volume 4595 of LNCS, pages 187–203. Springer, 2007.

\item[{[21]}] J. Barnat, L. Brim, I. Černá, and P. Šimeček. DiVinE – The Distributed Verification Environment. In M. Leucker and J. van de Pol, editors, Proceedings of 4th International Workshop on Parallel and Distributed Methods in verifiCation, pages 89–94, July 2005.

\item[{[22]}] J. Barnat, V. Forejt, M. Leucker, and M. Weber. DivSPIN – A SPIN compatible distributed model checker. In M. Leucker and J. van de Pol, editors, Proceedings of 4th International Workshop on Parallel and Distributed Methods in verifiCation, pages 95–100, July 2005.

\item[{[23]}] F. J. Barros. Dynamic structure discrete event system specification: a new formalism for dynamic structure modeling and simulation. In Proceedings of the 1995 Winter Simulation Multiconference, pages 781–785, 1995.

\item[{[24]}] F. J. Barros. Modeling formalisms for dynamic structure systems. ACM Transactions on Modeling and Computer Simulation, 7:501–515, 1997.

\item[{[25]}] F. J. Barros. Abstract simulators for the DSDE formalism. In Proceedings of the 1998 Winter simulation Multiconference, pages 407–412, 1998.

\item[{[26]}] F. Bause and P. S. Kritzinger. Stochastic petri nets - an introduction to the theory (2. ed.). In Advanced studies of computer science, 1996.

\item[{[27]}] G. Behrmann, A. David, K. G. Larsen, P. Pettersson, and W. Yi. Developing UPPAAL over 15 years. Softw., Pract. Exper., 41(2):133–142, 2011.

\item[{[28]}] L. A. Belady, M. W. Blasgen, C. J. Evangelisti, and R. D. Tennison. A computer graphics system for block diagram problems. IBM Systems Journal, 10(2):143–161, 1971.

\item[{[29]}] B. Berthomieu and F. Vernadat. Time petri nets analysis with TINA. In Third International Conference on the Quantitative Evaluation of Systems (QEST 2006), 11-14 September 2006, Riverside, California, USA, pages 123–124. IEEE Computer Society, 2006.

\item[{[30]}] Y. Bertot and P. Castéran. Interactive Theorem Proving and Program Development - Coq’Art: The Calculus of Inductive Constructions. Texts in Theoretical Computer Science. An EATCS Series. Springer, 2004.

\item[{[31]}] J. Bézivin and B. Jean. On the unification power of models. Vital And Health Statistics. Series 20 Data From The National Vitalstatistics System Vital Health Stat 20 Data Natl Vital Sta, 4(2):171–188, 2005.

\item[{[32]}] O. Biberstein, D. Buchs, and N. Guelfi. Object-Oriented Nets with Algebraic Specifications: The \{CO-OPN/2\} Formalism. In G. Agha, F. D. Cindio, and G. Rozenberg, editors, Advances in Petri Nets on Object-Orientation, LNCS, pages 70–127. Springer, 2001.

\item[{[33]}] A. Biere, A. Cimatti, E. M. Clarke, and Y. Zhu. Symbolic model checking without bdds. In TACAS ’99: Proceedings of the 5th International Conference on Tools and Algorithms for Construction and Analysis of Systems, pages 193–207, London, UK, 1999. Springer.

\item[{[34]}] A. Biere and K. Heljanko. Hardware model checking contest, 2013. \url{http://fmv.jku.at/hwmcc13/}.

\item[{[35]}] G. Birtwistle, O. J. Dahl, B. Myhrhaug, and K. Nygaard. Speed-Up in Chemical Engineering Design. Auerbach Publishers, Inc., 1973.

\item[{[36]}] T. Blochwitz, M. Otter, J. Åkesson, M. Arnold, C. Clauss, H. Elmqvist, M. Friedrich, A. Junghanns, J. Mauss, D. Neumerkel, H. Olsson, and A. Viel. Functional mockup interface 2.0: The standard for tool independent exchange of simulation models. In Proceedings of the 9th International Modelica Conference, pages 173–184. The Modelica Association, 2012.

\item[{[37]}] D. Blouin. Modeling Languages for Requirements Engineering and Quantitative Analysis of Embedded Systems. PhD thesis, University of South Britanny, France, 2013.

\item[{[38]}] D. Blouin and H. Giese. Combining Requirements, Use Case Maps and AADL Models for Safety-Critical Systems Design. In 2016 42th Euromicro Conference on Software Engineering and Advanced Applications (SEAA), pages 266–274, Aug 2016.

\item[{[39]}] E. L. Bóbeda, M. Colange, and D. Buchs. Building a symbolic model checker from formal language description. In ACSD, pages 50–59. IEEE Computer Society, 2015.

\item[{[40]}] B. Boehm. A spiral model of software development and enhancement. SIGSOFT Softw. Eng. Notes, 11(4):14–24, Aug. 1986.

\item[{[41]}] J.-S. Bolduc and H. Vangheluwe. Expressing ODE models as DEVS: Quantization approaches. In Proceedings of the AIS’2002 Conference (AI, Simulation and Planning in High Autonomy Systems), April 2002, Lisboa, Portugal/F. Barros and N. Giambiasi (eds.), pages 163–169, 2002.

\item[{[42]}] M. Bozzano, R. Cavada, A. Cimatti, J.-P. Katoen, V. Y. Nguyen, T. Noll, and X. Olive. Formal verification and validation of AADL models. Proc. ERTS, 2010.

\item[{[43]}] P. C. Breedveld. Physical Systems Theory in Terms of Bond Graphs. Phd thesis, University of Twente, 1984.

\item[{[44]}] P. C. Breedveld. Multibond-graph elements in physical systems theory. Journal of the Franklin Institute, 319:1–36, 1985.

\item[{[45]}] R. D. Brennan. Digital simulation for control system design. In Proceedings of the SHARE design automation project on - DAC ’66, DAC ’66, pages 1.1–1.12, New York, New York, USA, 1966. ACM Press.

\item[{[46]}] J. F. Broenink. Computer-Aided Physical-Systems Modeling and Simulation: A Bond-Graph Approach. Phd thesis, University of Twente, Enschede, The Netherlands, mar 1990.

\item[{[47]}] J. F. Broenink. Bond-graph modeling in modelica. pages 137–141. 9th European Simulation Symposium (ESS97), Passau, Germany, 1997.

\item[{[48]}] J. F. Broenink. Modelling, simulation and analysis with 20-sim. Journal A, 38:22 –25, sep 1997.

\item[{[49]}] J. F. Broenink. 20-sim software for hierarchical bond-graph/block-diagram models. Simulation Practice and Theory, nr. 7, pages 481–492, 1999.

\item[{[50]}] J. F. Broenink and C. Kleijn. Computer-aided design of mechatronic systems using 20-sim 3.0. In Proc WESIC, Workshop on Scientific and Industrial Collaboration, Newport, UK, 1999. IEE.

\item[{[51]}] D. Broman, C. Brooks, L. Greenberg, E. A. Lee, M. Masin, S. Tripakis, and M. Wetter. Determinate composition of fmus for co-simulation. In Proceedings of the Eleventh ACM International Conference on Embedded Software, EMSOFT ’13, pages 2:1–2:12, Piscataway, NJ, USA, 2013. IEEE Press.

\item[{[52]}] J. Brooks, F.P. No silver bullet essence and accidents of software engineering. Computer, 20(4):10 –19, april 1987.

\item[{[53]}] F. Brown. Direct application of the loop rule to bond graphs. Journal of Dynamic Systems, Measurement and Control, 94:253–261, 1972.

\item[{[54]}] R. Bruttomesso, M. Deters, and A. Griggio. Satisfiability modulo theories competition, 2014. \url{http://www.smtcomp.org/}.

\item[{[55]}] R. E. Bryant. Symbolic Boolean Manipulation with Ordered Binary-Decision Diagrams. ACM Comput. Surv., 24(3):293–318, 1992.

\item[{[56]}] D. Buchs, S. Hostettler, A. Marechal, and M. Risoldi. AlPiNA: A symbolic model checker. In J. Lilius and W. Penczek, editors, Applications and Theory of Petri Nets, volume 6128 of Lecture Notes in Computer Science, pages 287–296. Springer, 2010.

\item[{[57]}] D. Buchs, S. Hostettler, A. Marechal, and M. Risoldi. Alpina: An algebraic petri net analyzer. In J. Esparza and R. Majumdar, editors, Tools and Algorithms for the Construction and Analysis of Systems, volume 6015 of Lecture Notes in Computer Science, pages 349–352. Springer, 2010.

\item[{[58]}] J. R. Burch, E. M. Clarke, D. E. Long, K. L. McMillan, and D. L. Dill. Symbolic model checking for sequential circuit verification. IEEE Transactions on Computer-Aided Design of Integrated Circuits and Systems, 13:401–424, 1994.

\item[{[59]}] J. R. Burch, E. M. Clarke, K. L. McMillan, D. L. Dill, and L. J. Hwang. Symbolic model checking: $10^{20}$ states and beyond. Information and Computation, 98(2):142–170, 1992.

\item[{[60]}] R. L. Burden and J. D. Faires. Numerical Analysis. Cengage Learning, 9 edition, 2010.

\item[{[61]}] J. Butcher. Numerical Methods for Ordinary Differential Equations. John Wiley \& Sons, Ltd, Chichester, UK, jun 2003.

\item[{[62]}] L. Cabac, M. Haustermann, and D. Mosteller. Renew 2.5 - towards a comprehensive integrated development environment for petri net-based applications. In F. Kordon and D. Moldt, editors, Application and Theory of Petri Nets and Concurrency - 37th International Conference, PETRI NETS 2016, Toruń, Poland, June 19-24, 2016. Proceedings, volume 9698 of Lecture Notes in Computer Science, pages 101–112. Springer, 2016.

\item[{[63]}] Case Study AADL Models. \url{https://gitlab.telecom-paristech.fr/mbe-tools/RAMSES/tree/develop/examples/line_follower_mindstorm_robot}.

\item[{[64]}] F. Casella. Proceedings of the 7th International Modelica Conference, Como, Italy. 2009.

\item[{[65]}] F. Celier. Combined Continuous/Discrete System Simulation by Use of Digital Computers. Techniques and Tools. ETH No 6483 Swiss Federal Institute of Technology Zurich, Switzerland, 1979.

\item[{[66]}] F. Cellier, P. Fritzson, E. Lee, and D. Broman. Proceedings of the 4th International Workshop on Equation-Based Object-Oriented Languages and Tools, (EOOLT’2011), Zürich, Switzerland. Linköping University Electronic Press, 2011.

\item[{[67]}] F. E. Cellier. Continuous System Modeling. Springer Verlag, 1991.

\item[{[68]}] F. E. Cellier and E. Kofman. Continuous System Simulation. Springer-Verlag, Berlin, Heidelberg, 2006.

\item[{[69]}] F. E. Cellier and E. Kofman. Continuous System Simulation. Springer Science \& Business Media, 2006.

\item[{[70]}] B. Chen and H. Vangheluwe. Symbolic flattening of DEVS models. In Proceedings of the 2010 Summer Simulation Multiconference, pages 209–218, 2010.

\item[{[71]}] G. Chiola. Characterization of timed well-formed petri nets behavior by means of occurrence equations. In Proceedings of the Sixth International Workshop on Petri Nets and Performance Models, PNPM 1995, Durham, NC, USA, October 3-6, 1995, pages 127–136. IEEE Computer Society, 1995.

\item[{[72]}] A. C. H. Chow and B. P. Zeigler. Parallel DEVS: a parallel, hierarchical, modular, modeling formalism. In Proceedings of the 1994 Winter Simulation Multiconference, pages 716–722, 1994.

\item[{[73]}] A. C. H. Chow, B. P. Zeigler, and D. H. Kim. Abstract simulator for the parallel DEVS formalism. In AI, Simulation, and Planning in High Autonomy Systems, pages 157–163, 1994.

\item[{[74]}] E. M. Clarke, Jr., O. Grumberg, and D. A. Peled. Model checking. MIT Press, Cambridge, MA, USA, 1999.

\item[{[75]}] C. Clauss. Proceedings of the 8th International Modelica Conference, Dresden, Germany. 2011.

\item[{[76]}] N. S. E. D. M. Committee et al. Final report of the Model Based Engineering (MBE) subcommittee. National Defense Industrial Association (NDIA), 2011.

\item[{[77]}] O. Coudert, C. Berthet, and J. C. Madre. Formal boolean manipulations for the verification of sequential machines. In EURO-DAC ’90: Proceedings of the conference on European design automation, pages 57–61, Los Alamitos, CA, USA, 1990. IEEE Computer Society Press.

\item[{[78]}] B. Curtis, M. I. Kellner, and J. Over. Process Modeling. Communications of the ACM, 35(9):75–90, 1992.

\item[{[79]}] I. Dávid, H. Vangheluwe, and Y. Van Tendeloo. Translating engineering workflow models to DEVS for performance evaluation. In Proceedings of the 2018 Winter Simulation Conference, pages 616–627. IEEE Press, 2018.

\item[{[80]}] B. Davis, H. Porta, and J. Uhl. Calculus \& Mathematica Vector Calculus: Measuring in Two and Three Dimensions. Addison-Wesley, 1994.

\item[{[81]}] J. De Lara and H. Vangheluwe. Computer Aided Multi-paradigm Modelling to Process Petri-nets and Statecharts. In International Conference on Graph Transformation, pages 239–253. Springer, 2002.

\item[{[82]}] J. Delange, P. Feiler, and E. Neil. Incremental Life Cycle Assurance of Safety-Critical Systems. In 8th European Congress on Embedded Real Time Software and Systems (ERTS 2016), Toulouse, France, Jan. 2016.

\item[{[83]}] J. Denil, B. Meyers, P. D. Meulenaere, and H. Vangheluwe. Explicit semantic adaptation of hybrid formalisms for FMI co-simulation. In Proceedings of the Symposium on Theory of Modeling \& Simulation: DEVS Integrative M\&S Symposium, part of the 2015 Spring Simulation Multiconference, SpringSim ’15, Alexandria, VA, USA, April 12-15, 2015, pages 99–106, 2015.

\item[{[84]}] P. Derler, E. A. Lee, M. Törngren, and S. Tripakis. Cyber-physical system design contracts. In 2013 ACM/IEEE International Conference on Cyber-Physical Systems (ICCPS), pages 109–118, April 2013.

\item[{[85]}] M. Diaz. Petri Nets: Fundamental Models, Verification and Applications. ISTE. Wiley, 2013.

\item[{[86]}] J. Dixhoorn, van. Bond graphs and the challenge of a unified modelling theory of physical systems. In F. E. Cellier, editor, Progress in Modelling and Simulation, pages 207 – 245. Academic Press, 1982.

\item[{[87]}] H. Ehrig and B. Mahr. Fundamentals of Algebraic Specification 1: Equations und Initial Semantics, volume 6 of EATCS Monographs on Theoretical Computer Science. Springer, 1985.

\item[{[88]}] H. Elmqvist. A Structured Model Language for Large Continuous Systems. PhD thesis, 1978.

\item[{[89]}] H. Elmqvist, D. Bruck, and M. Otter. Dymola—User’s Manual. Dynasim AB, Research Park Ideon, SE-223 70, Lund, Sweden, 1996.

\item[{[90]}] H. Elmqvist, S.-E. Mattsson, and M. Otter. Modelica—a language for physical system modeling, visualization and interaction. In Proceedings of the 1999 IEEE Symposium on Computer-Aided Control System Design (CACSD’99), Hawaii, 1999.

\item[{[91]}] H. Elmqvist, M. Otter, S.-E. Mattsson, and H. Olsson. Modeling, Simulation, and Optimization with Modelica and Dymola,Book draft. Dynasim AB, Research Park Ideon, SE-223 70, Lund, Sweden, 2002.

\item[{[92]}] E. A. Emerson, S. Jha, and D. Peled. Combining partial order and symmetry reductions. In TACAS, pages 19–34, 1997.

\item[{[93]}] T. Ernst, S. Jahnichen, M. Klose, and R. Chaussee. The architecture of the smile/m simulation environment. In Proceedings 15th IMACS World Congress on Scientific Computation, Modelling and Applied Mathematics, 01 1998.

\item[{[94]}] S. Esmaeilsabzali, N. A. Day, J. M. Atlee, and J. Niu. Deconstructing the semantics of big-step modelling languages. Requirements Engineering, 15(2):235–265, 2010.

\item[{[95]}] J. A. Estefan. Survey of Model-Based Systems Engineering (MBSE) Methodologies. Technical Report Rev. B, INCOSE MBSE Initiative, May 2008.

\item[{[96]}] S. Evangelista and J.-F. Pradat-Peyre. Memory efficient state space storage in explicit software model checking. In in Proceedings of the 12th International SPIN Workshop on Model Checking of Software, volume 3639, pages 43–57. Springer, 2005.

\item[{[97]}] J. Farkas. Theorie der einfachen ungleichungen. Journal für die reine und angewandte Mathematik, 124:1–27, 1902.

\item[{[98]}] P. Feiler, J. Hansson, D. de Niz, and L. Wrage. System Architecture Virtual Integration: An Industrial Case Study. Technical Report CMU/SEI-2009-TR-017, Software Engineering Institute, Carnegie Mellon University, Pittsburgh, PA, 2009.

\item[{[99]}] P. H. Feiler. Modeling of System Families. Technical report, Software Engineering Institute (SEI), 2007.

\item[{[100]}] P. H. Feiler and D. P. Gluch. Model-Based Engineering with AADL: An Introduction to the SAE Architecture Analysis \& Design Language. Addison-Wesley Professional, 1st edition, 2012.

\item[{[101]}] Design and Analysis of Cyber-Physical Systems: AADL and Avionics Systems. \url{https://resources.sei.cmu.edu/library/asset-view.cfm?assetid=48553}.

\item[{[102]}] M. Felleisen, R. Findler, M. Flatt, and S. Krishnamurthi. The drscheme project: An overview. In Proceedings of the ACM SIGPLAN 1998 Conference on Programming Language Design and Implementation (PLDI’98), Montreal, Canada. ACM, 1998.

\item[{[103]}] A. Fernström, I. Axelsson, P. Fritzson, A. Sandholm, and A. Pop. OMNotebook – Interactive WYSIWYG Book Software for Teaching Programming. Proc. of the Workshop on Developing Computer Science Education – How Can It Be Done? Linköping University, Dept. Computer \& Inf. Science, Linköping, Sweden, 2006.

\item[{[104]}] P. Fritzson. Proceedings of SIMS’99—The 1999 Conference of the Scandinavian Simulation Society. 1999.

\item[{[105]}] P. Fritzson. Proceedings of Modelica 2000 Workshop. Lund University, Lund, Sweden, 2000.

\item[{[106]}] P. Fritzson. Proceedings of the 3rd International Modelica Conference. Linköping University, Linköping, Sweden, 2003.

\item[{[107]}] P. Fritzson. Principles of Object Oriented Modeling and Simulation with Modelica 2.1. Wiley-IEEE Press, 2004.

\item[{[108]}] P. Fritzson. Proceedings of the 2nd Conference on Modeling and Simulation for Safety and Security (SimSafe’2005). Linköping University, Linköping, Sweden, 2005.

\item[{[109]}] P. Fritzson. Proceedings of the 4rd International Conference on Modeling and Simulation in Biology, Medicine, and Biomedical Engineering (BioMedSim’2005). Linköping University, Linköping, Sweden, 2005.

\item[{[110]}] P. Fritzson. Mathmodelica - an object oriented mathematical modeling and simulation environment. 10(1), 2006.

\item[{[111]}] P. Fritzson. Introduction to Modeling and Simulation of Technical and Physical Systems with Modelica. Wiley-IEEE Press, 2011.

\item[{[112]}] P. Fritzson. Principles of Object Oriented Modeling and Simulation with Modelica 3.3: A Cyber-Physical Approach. Wiley IEEE Press, 2014.

\item[{[113]}] P. Fritzson and K.-F. Berggren. A pseudopotential calculation of the density of states of expanded crystalline mercury. Solid State Communications, 19(4):385 – 387, 1976.

\item[{[114]}] P. Fritzson and P. Bunus. Modelica—a general object-oriented language for continuous and discrete-event system modeling and simulation. In Proceedings of the 35th Annual Simulation Symposium, San Diego, California. IEEE, 2002.

\item[{[115]}] P. Fritzson, F. Cellier, and D. Broman. Proceedings of the 2nd International Workshop on Equation-Based Object-Oriented Languages and Tools, (EOOLT’2008), in conjunction with ECOOP, Paphos, Cyprus. Linköping University Electronic Press, 2008.

\item[{[116]}] P. Fritzson, F. Cellier, and C. Nytsch Geusen. Proceedings of the 1st International Workshop on Equation-Based Object-Oriented Languages and Tools, (EOOLT’2007), in conjunction with ECOOP, Berlin. Linköping University Electronic Press, 2007.

\item[{[117]}] P. Fritzson, V. Engelson, and J. Gunnarsson. An integrated modelica environment for modeling, documentation and simulation. In Proceedings of Summer Computer Simulation Conference ’98, Reno, Nevada, USA, 1998.

\item[{[118]}] P. Fritzson, E. Lee, F. Cellier, and D. Broman. Proceedings of the 3rd International Workshop on Equation-Based Object-Oriented Languages and Tools, (EOOLT’2010), in conjunction with MODELS, Oslo, Norway. Linköping University Electronic Press, 2010.

\item[{[119]}] P. Fritzson, L. Viklund, J. Herber, and D. Fritzson. Industrial application of object-oriented mathematical modeling and computer algebra in mechanical analysis. 1992.

\item[{[120]}] P. Fritzson, L. Viklund, O. J. Herber, and D. Fritzson. High-level mathematical modeling and programming. IEEE Software, 12(4):77–87, July 1995.

\item[{[121]}] H. Garavel, R. Mateescu, and I. Smarandache. Parallel state space construction for model-checking. In Proceedings of the 8th international SPIN workshop on Model checking of software, SPIN ’01, pages 217–234, New York, NY, USA, 2001. Springer-Verlag New York, Inc.

\item[{[122]}] W. Gellert and V. N. Reinhold. The VNR concise encyclopedia of mathematics / W. Gellert ... [et al.], editors ; K.A. Hirsch, H. Reichardt, scientific advisors. Van Nostrand Reinhold New York, 2nd ed. edition, 1989.

\item[{[123]}] D. Gianni, A. D’Ambrogio, and A. Tolk. Introduction to the modeling and simulation-based systems engineering handbook., 2014.

\item[{[124]}] H. Giese, S. Neumann, O. Niggemann, and B. Schätz. Model-based integration. In Proceedings of the 2007 International Dagstuhl conference on Model-based engineering of embedded real-time systems, pages 17–54. Springer-Verlag, 2007.

\item[{[125]}] P. Godefroid. Using partial orders to improve automatic verification methods. In E. Clarke and R. Kurshan, editors, Computer-Aided Verification, volume 531 of Lecture Notes in Computer Science, pages 176–185. Springer, 1991.

\item[{[126]}] C. Gomes, C. Thule, D. Broman, P. G. Larsen, and H. Vangheluwe. Co-simulation: State of the art. CoRR, abs/1702.00686, 2017.

\item[{[127]}] R. Goud, K. M. van Hee, R. D. J. Post, and J. M. E. M. van der Werf. Petriweb: A repository for petri nets. In S. Donatelli and P. S. Thiagarajan, editors, Petri Nets and Other Models of Concurrency - ICATPN 2006, 27th International Conference on Applications and Theory of Petri Nets and Other Models of Concurrency, Turku, Finland, June 26-30, 2006, Proceedings, volume 4024 of Lecture Notes in Computer Science, pages 411–420. Springer, 2006.

\item[{[128]}] F. Gracer and R. A. Myers. Graphic Computer-assisted Design of Optical Filters. IBM Journal of Research and Development, 13(2):172–178, mar 1969.

\item[{[129]}] J. Granda. Computer generation of physical system differential equations using bond graphs. Journal of the Franklin Institute, 319:243–256, 1985.

\item[{[130]}] C. Greer, M. Burns, D. Wollman, and E. Griffor. Cyber-physical systems and internet of things. NIST Special Publication 1900-202.

\item[{[131]}] E. Guerra, J. De Lara, D. S. Kolovos, R. F. Paige, and O. M. Dos Santos. transML: A family of languages to model model transformations. In International Conference on Model Driven Engineering Languages and Systems, pages 106–120. Springer, 2010.

\item[{[132]}] A. Hamez, L. Hillah, F. Kordon, A. Linard, E. Paviot-Adet, X. Renault, and Y. Thierry-Mieg. New features in CPN-AMI 3: focusing on the analysis of complex distributed systems. In Sixth International Conference on Application of Concurrency to System Design (ACSD 2006), 28-30 June 2006, Turku, Finland, pages 273–275. IEEE Computer Society, 2006.

\item[{[133]}] J. Hansson, S. Helton, and P. Feiler. ROI Analysis of the System Architecture Virtual Integration Initiative. Technical Report CMU/SEI-2018-TR-002, Software Engineering Institute, Carnegie Mellon University, Pittsburgh, PA, 2018.

\item[{[134]}] D. Harel. Statecharts: A visual formalism for complex systems. Science of Computer Programming, 8(3):231–274, 1987.

\item[{[135]}] D. Harel. Statecharts: A Visual Formulation for Complex Systems. Sci. Comput. Program., 8(3):231–274, 1987.

\item[{[136]}] D. Harel. On visual formalisms. Commun. ACM, 31(5):514–530, May 1988.

\item[{[137]}] D. Harel and E. Gery. Executable object modeling with statecharts. IEEE Computer, 30(7):31–42, 1997.

\item[{[138]}] D. Harel and H. Kugler. The rhapsody semantics of statecharts (or, on the executable core of the UML). In Integration of Software Specification Techniques for Applications in Engineering, volume 3147 of Lecture Notes in Computer Science, pages 325–354. Springer Berlin Heidelberg, 2004.

\item[{[139]}] D. Harel and H. Kugler. The Rhapsody semantics of Statecharts (or, on the executable core of the UML). In H. Ehrig, W. Damm, J. Desel, M. Große-Rhode, W. Reif, E. Schnieder, and E. Westkämper, editors, Integration of Software Specification Techniques for Applications in Engineering, pages 325–354. Springer, Berlin Heidelberg, 2004.

\item[{[140]}] D. Harel and A. Naamad. The STATEMATE semantics of Statecharts. ACM Transactions on Software Engineering Methodology, 5(4):293–333, 1996.

\item[{[141]}] D. Harel and A. Naamad. The STATEMATE Semantics of Statecharts. ACM Transactions on Software Engineering and Methodology, 5(4):293–333, 1996.

\item[{[142]}] R. Hebig, A. Seibel, and H. Giese. On the Unification of Megamodels. In Electronic Communications of the EASST: Proceedings of the 4th International Workshop on Multi-Paradigm Modeling, volume 42, 2011.

\item[{[143]}] F. J. Heemstra and R. J. Kusters. Dealing with risk: a practical approach. Journal of Information Technology, 11(4):333–346, 1996.

\item[{[144]}] D. C. Heggie. J. fauvel, r. flood, m. shortland and r. wilson, let newton be! (clarendon press, oxford, 1989). Proceedings of the Edinburgh Mathematical Society, 33(1):166–168, 1990.

\item[{[145]}] F. Heidenreich, J. Kopcsek, and U. Aßmann. Safe Composition of Transformations. In International Conference on Theory and Practice of Model Transformations, pages 108–122. Springer, 2010.

\item[{[146]}] L. Hillah and F. Kordon. Petri nets repository: A tool to benchmark and debug petri net tools. In W. M. P. van der Aalst and E. Best, editors, Application and Theory of Petri Nets and Concurrency - 38th International Conference, PETRI NETS 2017, Zaragoza, Spain, June 25-30, 2017, Proceedings, volume 10258 of Lecture Notes in Computer Science, pages 125–135. Springer, 2017.

\item[{[147]}] L. Hillah, F. Kordon, L. Petrucci, and N. Trèves. PNML framework: An extendable reference implementation of the petri net markup language. In J. Lilius and W. Penczek, editors, Applications and Theory of Petri Nets, 31st International Conference, PETRI NETS 2010, Braga, Portugal, June 21-25, 2010. Proceedings, volume 6128 of Lecture Notes in Computer Science, pages 318–327. Springer, 2010.

\item[{[148]}] J. Himmelspach and A. M. Uhrmacher. Sequential processing of PDEVS models. In Proceedings of the 3rd European Modeling \& Simulation Symposium, pages 239–244, 2006.

\item[{[149]}] G. J. Holzmann. The SPIN Model Checker. Addison Wesley, 2003.

\item[{[150]}] J. E. Hopcroft, R. Motwani, and J. D. Ullman. Introduction to Automata Theory, Languages, and Computation. Addison-Wesley Longman Publishing Co., Inc., Boston, MA, USA, 3r d edition, 2006.

\item[{[151]}] S. Hostettler, A. Marechal, A. Linard, M. Risoldi, and D. Buchs. High-Level Petri Net Model Checking with AlPiNA. Fundamenta Informaticae, 113, 2011.

\item[{[152]}] M.-H. Hwang. Generating finite-state global behavior of reconfigurable automation systems: DEVS approach. In Proceedings of the International Conference on Automation Science and Engineering, pages 254 – 260, 2005.

\item[{[153]}] INCOSE. Systems Engineering Vision 2020. Technical Report INCOSE-TP-2004-004-02, International Council on Systems Engineering (INCOSE), 2007.

\item[{[154]}] T. INCOSE. Systems engineering vision 2020. INCOSE, San Diego, CA, accessed Jan, 26:2019, 2007.

\item[{[155]}] ITU URN Specification. \url{http://www.itu.int/rec/T-REC-Z.150/en/}.

\item[{[156]}] iTunes Crashes when Ripping. \url{https://discussions.apple.com/thread/247969}.

\item[{[157]}] J. F. Jensen, T. Nielsen, L. K. Oestergaard, and J. Srba. TAPAAL and reachability analysis of P/T nets. T. Petri Nets and Other Models of Concurrency, 11:307–318, 2016.

\item[{[158]}] K. Jensen. Coloured Petri Nets - Basic Concepts, Analysis Methods and Practical Use - Volume 1, Second Edition. Monographs in Theoretical Computer Science. An EATCS Series. Springer, 1996.

\item[{[159]}] M. Järvisalo, D. Le Berre, and O. Roussel. Sat competitions, 2011. \url{http://www.satcompetition.org/}.

\item[{[160]}] D. C. Karnopp, D. L. Margolis, and R. C. Rosenberg. System Dynamics: Modeling and Simulation of Mechatronic Systems. Wiley, 2006.

\item[{[161]}] D. C. Karnopp and R. C. Rosenberg. Analysis and Simulation of Multiport Systems - The Bond Graph Approach to Physical System Dynamics. MIT Press, 1968.

\item[{[162]}] D. C. Karnopp and R. C. Rosenberg. System Dynamics, a Unified Approach. Wiley, 1975.

\item[{[163]}] S. Katz and D. Peled. An efficient verification method for parallel and distributed programs. In Linear Time, Branching Time and Partial Order in Logics and Models for Concurrency, School/Workshop, pages 489–507, London, UK, 1989. Springer.

\item[{[164]}] E. Kindler. The epnk: An extensible petri net tool for PNML. In L. M. Kristensen and L. Petrucci, editors, Applications and Theory of Petri Nets - 32nd International Conference, PETRI NETS 2011, Newcastle, UK, June 20-24, 2011. Proceedings, volume 6709 of Lecture Notes in Computer Science, pages 318–327. Springer, 2011.

\item[{[165]}] V. Klebanov, P. Müller, N. Shankar, G. T. Leavens, V. Wüstholz, E. Alkassar, R. Arthan, D. Bronish, R. Chapman, E. Cohen, M. Hillebrand, B. Jacobs, K. R. M. Leino, R. Monahan, F. Piessens, N. Polikarpova, T. Ridge, J. Smans, S. Tobies, T. Tuerk, M. Ulbrich, and B. Weiß. The 1st Verified Software Competition: Experience report. In M. Butler and W. Schulte, editors, FM’2011: 17th International Symposium on Formal Methods, volume 6664 of LNCS. Springer, 2011. \url{http://www.vscomp.org}.

\item[{[166]}] A. G. Kleppe, J. Warmer, J. B. Warmer, and W. Bast. MDA explained: the model driven architecture: practice and promise. Addison-Wesley Professional, 2003.

\item[{[167]}] S. Klikovits, R. Al-Ali, M. Amrani, A. Barisic, F. Barros, D. Blouin, E. Borde, D. Buchs, H. Giese, M. Goulao, M. Iacono, F. Leon, E. Navarro, P. Pelliccione, and K. Vanherpen. COST IC1404 WG1 Deliverable WG1.1: State-of-the- art on Current Formalisms used in Cyber-Physical Systems Development, Jan. 2019.

\item[{[168]}] C. Knieke and U. Goltz. An Executable Semantics for UML 2 Activity Diagrams. In Proceedings of the International Workshop on Formalization of Modeling Languages, page 3. ACM, 2010.

\item[{[169]}] F. Kordon and D. Buchs. Model Checking Contest @ Petri Nets, 2014. \url{http://mcc.lip6.fr}.

\item[{[170]}] F. Kordon, H. Garavel, L. Hillah, E. Paviot-Adet, L. Jezequel, C. Rodríguez, and F. Hulin-Hubard. Mcc’2015 - the fifth model checking contest. T. Petri Nets and Other Models of Concurrency, 11:262– 273, 2016.

\item[{[171]}] F. Kordon, A. Linard, D. Buchs, M. Colange, S. Evangelista, J.-F. Jensen, K. Lampka, N. Lohmann, E. Paviot-Adet, Y. Thierry-Mieg, and H. Wimmel. Report on the Model Checking Contest at Petri Nets 2011. Transactions on Petri Nets and Other Models of Concurrency, 2011.

\item[{[172]}] C. Kral and A. Haumer. Proceedings of the 6th International Modelica Conference, Vienna, Austria. 2006.

\item[{[173]}] G. Kron. Diakoptics: The Piecewise Solution of Large-scale Systems. McDonald \& Co. Ltd., 1963.

\item[{[174]}] P. Kruchten. The Rational Unified Process An Introduction. NY, 1980.

\item[{[175]}] T. S. Kuhn, O. Neurath, J. Dewey, and T. S. Kuhn. The Structure of scientific revolutions. Number ed.-in-chief: Otto Neurath ; Vol. 2 No. 2 in International encyclopedia of unified science Foundations of the unity of science. Chicago Univ. Press, Chicago, Ill, 2nd ed., enlarged edition, 1994. OCLC: 258260085.

\item[{[176]}] T. Kühne. Matters of (meta-)modeling. Software and System Modeling, 5(4):369–385, 2006.

\item[{[177]}] Y. W. Kwon, H. C. Park, S. H. Jung, and T. G. Kim. Fuzzy-DEVS formalism: concepts, realization and application. In Proceedings of AI, Simulation and Planning in High Autonomy Systems, pages 227 – 234, 1996.

\item[{[178]}] J. d. Lara, H. Vangheluwe, and P. J. Mosterman. Modelling and analysis of traffic networks based on graph transformation. Formal Methods for Automation and Safety in Railway and Automotive Systems, 11, 2004.

\item[{[179]}] B. R. Larson, P. Chalin, and J. Hatcliff. BLESS: Formal Specification and Verification of Behaviors for Embedded Systems with Software. In NASA Formal Methods, 5th International Symposium, NFM 2013, Moffett Field, CA, USA, May 14-16, 2013. Proceedings, pages 276–290, 2013.

\item[{[180]}] D. P. Y. Lawrence, C. Gomes, J. Denil, H. Vangheluwe, and D. Buchs. Coupling petri nets with deterministic formalisms using co-simulation. In Proceedings of the Symposium on Theory of Modeling \& Simulation, TMS/DEVS 2016, part of the 2016 Spring Simulation Multiconference, SpringSim ’16, Pasadena, CA, USA, April 3-6, 2016, page 6, 2016.

\item[{[181]}] E. A. Lee. The problem with threads. Computer, 39(5):33–42, 2006.

\item[{[182]}] Lego-Mindstorm Website. \url{https://www.lego.com/mindstorms/}.

\item[{[183]}] D. Lempia and S. Miller. Requirements Engineering Management Findings Report. Technical report, Federal Aviation Administration, 2009.

\item[{[184]}] D. Lempia and S. Miller. Requirements Engineering Management Handbook. Technical report, Federal Aviation Administration, 2009.

\item[{[185]}] E.-L. Lengquist Sandelin and S. Monemar. DrModelica–An Experimental Computer-Based Teaching Material for Modelica. Master’s thesis, LITH-IDA-Ex03/3, Department of Computer and Information Science, Linköping University, Linköping, Sweden, 2003.

\item[{[186]}] E.-L. Lengquist-Sandelin, S. Monemar, P. Fritzson, and P. Bunus. Drmodelica - an interactive tutoring environment for modelica. In Proceedings of the 3rd International Modelica Conference, Nov. 3-4, Linköping, Sweden, 2003.

\item[{[187]}] C. Lin, A. Chandhury, A. B. Whinston, and D. C. Marinescu. Logical inference of horn clauses in petri net models. In CSD-TR-91-031. Purdue University, 1991.

\item[{[188]}] G. Liu and K. Barkaoui. A survey of siphons in Petri nets. Information Sciences, 363:198–220, 10 2016.

\item[{[189]}] S. Liu, V. Stavridou, B. Dutertre, and S. L. V. Stavridou. The Practice of Formal Methods in Safety Critical Systems. Journal of Systems and Software, 28:28–77, 1995.

\item[{[190]}] E. López Bóbeda, M. Colange, and D. Buchs. StrataGEM: A Generic Petri Net Verification Framework. In G. Ciardo and E. Kindler, editors, Proceedings of the 35th International Conference on Application and Theory of Petri Nets and Concurrency, June 2014.

\item[{[191]}] F. Lorenz. Modelling System 1, Users Manual. iege, Belgium, 1997.

\item[{[192]}] LRDE. Spot home page. \url{http://spot.lip6.fr/}.

\item[{[193]}] L. Lúcio, S. bin Abid, S. Rahman, V. Aravantinos, R. Kuestner, and E. Harwardt. Process-Aware Model-driven Development Environments. In MODELS (Satellite Events), pages 405–411, 2017.

\item[{[194]}] L. Lúcio, S. Mustafiz, J. Denil, B. Meyers, and H. Vangheluwe. The Formalism Transformation Graph as a Guide to Model Driven Engineering. SOCS-TR-2012.1, McGill University, 2012.

\item[{[195]}] L. Lúcio, S. Mustafiz, J. Denil, H. Vangheluwe, and M. Jukss. Ftg+pm: An integrated framework for investigating model transformation chains. In F. Khendek, M. Toeroe, A. Gherbi, and R. Reed, editors, Proceedings of SDL 2013: Model-Driven Dependability Engineering: 16th International SDL Forum, pages 182–202, Berlin, Heidelberg, 2013. Springer.

\item[{[196]}] L. Lúcio, S. Mustafiz, J. Denil, H. Vangheluwe, and M. Jukss. FTG+PM: An Integrated Framework for Investigating Model Transformation Chains. In SDL 2013: Model-Driven Dependability Engineering, volume 7916 of Lecture Notes in Computer Science, pages 182–202, Berlin, Heidelberg, 2013. Springer Berlin Heidelberg.

\item[{[197]}] A. Marechal and D. Buchs. Generalizing the Compositions of Petri Nets Modules. Fundamenta Informaticae, 137(1), 2015.

\item[{[198]}] R. Margaret. Capability Maturity Model (CMM). \url{https://searchsoftwarequality.techtarget.com/definition/Capability-Maturity-Model}. Last access: April 2019.

\item[{[199]}] D. C. Marinescu, M. Beaven, and R. Stansifer. A parallel algorithm for computing invariants of petri net models. In Petri Nets and Performance Models, 1991. PNPM91., Proceedings of the Fourth International Workshop on, pages 136–143. IEEE, 1991.

\item[{[200]}] MASIW Homepage. \url{https://forge.ispras.ru/projects/masiw-oss}.

\item[{[201]}] Mathworks. Simulink User’s Guide. MathWorks Inc., 2001.

\item[{[202]}] Mathworks. MATLAB User’s Guide. MathWorks Inc., 2002.

\item[{[203]}] S. Mattsson, M. Andersson, and K. Åström. Object-Oriented Modelling and Simulation. Marcel Dekker, 1993.

\item[{[204]}] S. E. Mattsson, H. Elmqvist, and J. F. Broenink. Modelica: The next generation modeling language - an international design effort. In Proceedings of First World Congress of System Simulation, Singapore, 1997.

\item[{[205]}] J. D. McGregor, D. P. Gluch, and P. H. Feiler. Analysis and Design of Safety-critical, Cyber-Physical Systems. Ada Lett., 36(2):31–38, May 2017.

\item[{[206]}] W. Michael and L. M. Kristensen. The Access-CPN Framework: A Tool for Interacting with the CPN Tools Simulator. In F. Giuliana and K. Wolf, editors, Applications and Theory of Petri Nets, pages 313–322, Berlin, Heidelberg, 2009. Springer Berlin Heidelberg.

\item[{[207]}] R. Milner. Communication and Concurrency. Prentice-Hall, Inc., Upper Saddle River, NJ, USA, 1989.

\item[{[208]}] A PID Controller for Lego Mindstorms Robots. \url{http://www.inpharmix.com/jps/PID_Controller_For_Lego_Mindstorms_Robots.html}.

\item[{[209]}] E. Mitchell and J. S.Gauthier. ACSL: Advanced Continuous Simulation Language-User Guide and Reference Manual. Mitchell \& Gauthier Assoc., Concord, Mass, 1986.

\item[{[210]}] T.-T. Mohsen, J. Remala, M. Sjolund, A. Pop, and P. Fritzson. Omsketch - graphical sketching in the openmodelica interactive book, omnotebook. In Proceedings of the 52th Scandinavian Conference on Simulation and Modeling (SIMS 2011), Vasteras, Sweden, 2011.

\item[{[211]}] P. J. Mosterman and G. Biswas. A theory of discontinuities in physical system models. Journal of the Franklin Institute, 335(3):401–439, apr 1998.

\item[{[212]}] P. J. Mosterman and H. Vangheluwe. Computer Automated Multi-Paradigm Modeling: An Introduction. Simulation, 80(9):433–450, sep 2004.

\item[{[213]}] P. J. Mosterman and H. Vangheluwe. Computer automated multi-paradigm modeling: An introduction. Simulation, 80(9):433–450, 2004.

\item[{[214]}] T. Murata, B. Shenker, and S. M. Shatz. Detection of ada static deadlocks using petri net invariants. IEEE Transactions on Software engineering, 15(3):314–326, 1989.

\item[{[215]}] T. Murata and D. Zhang. A predicate-transition net model for parallel interpretation of logic programs. IEEE Transactions on Software Engineering, 14(4):481–497, 1988.

\item[{[216]}] S. Mustafiz, J. Denil, L. Lúcio, and H. Vangheluwe. The FTG+PM framework for multi-paradigm modelling: An automotive case study. In Proceedings of the 6th International Workshop on Multi-Paradigm Modeling, pages 13–18. ACM, 2012.

\item[{[217]}] H. Nilsson. Proceedings of the 5th International Workshop on Equation-Based Object-Oriented Languages and Tools, (EOOLT’2013), Nottingham. Linköping University Electronic Press, 2013.

\item[{[218]}] J. J. Nutaro. Building Software for Simulation: Theory and Algorithms, with Applications in C++. Wiley, 1st edition, 2010.

\item[{[219]}] NXT-OSEK Homepage. \url{http://lejos-osek.sourceforge.net/}.

\item[{[220]}] NXT Programs Website. \url{http://www.nxtprograms.com/line_follower/steps.html}.

\item[{[221]}] Object Management Group. Precise Semantics of UML State Machines Specification, May 2019.

\item[{[222]}] Ocarina Website. \url{http://www.openaadl.org/ocarina.html}.

\item[{[223]}] P. C. Ölveczky, A. Boronat, and J. Meseguer. Formal Semantics and Analysis of Behavioral AADL Models in Real-Time Maude. In Formal Techniques for Distributed Systems, Joint 12th IFIP WG 6.1 International Conference, FMOODS 2010 and 30th IFIP WG 6.1 International Conference, FORTE 2010, Amsterdam, The Netherlands, June 7-9, 2010. Proceedings, pages 47–62, 2010.

\item[{[224]}] Open Source AADL Tool Environment (OSATE). \url{http://osate.org/}.

\item[{[225]}] M. Otter. Objektorientierte modellierung physikalischer systeme, teil 1: Übersicht. at Automatisierungstechnik, 47(1):A1–A4, 1999. LIDO-Berichtsjahr=1999,.

\item[{[226]}] M. Otter. Proceedings of the 2nd International Modelica Conference,Oberpfaffenhofen, Germany. 2002.

\item[{[227]}] Papyrus-RT Homepage. \url{https://www.eclipse.org/papyrus-rt/}.

\item[{[228]}] H. M. Paynter. Analysis and Design of Engineering Systems. MIT Press, 1961.

\item[{[229]}] F. Peter, M. Jirstrand, and J. Gunnarsson. Mathmodelica—an extensible modeling and simulation environment with integrated graphics and literate programming. In Proceedings of the 2nd International Modelica Conference, Oberpfaffenhofen, Germany, 2002.

\item[{[230]}] F. Pommereau. SNAKES: A flexible high-level petri nets library (tool paper). In R. R. Devillers and A. Valmari, editors, Application and Theory of Petri Nets and Concurrency - 36th International Conference, PETRI NETS 2015, Brussels, Belgium, June 21-26, 2015, Proceedings, volume 9115 of Lecture Notes in Computer Science, pages 254–265. Springer, 2015.

\item[{[231]}] L. Popova-Zeugmann. Time and Petri Nets. Springer, 2013.

\item[{[232]}] E. Posse, J. de Lara, and H. Vangheluwe. Processing causal block diagrams with graphgrammars in atom3. In Workshop on Applied Graph Transformation (AGT), pages 23–34, Grenoble, France, apr 2002. Springer, Berlin, Heidelberg.

\item[{[233]}] A. Pritsker and B. Alan. The GASP IV simulation language [by] A. Alan B. Pritsker. Wiley New York, 1974.

\item[{[234]}] S. Rahmoun, E. Borde, and L. Pautet. multi-objectives refinement of aadl models for the synthesis embedded systems (mu-ramses). In 2015 20th International Conference on Engineering of Complex Computer Systems (ICECCS).

\item[{[235]}] RAMSES Homepage. \url{https://mem4csd.telecom-paristech.fr/blog/index.php/ramses/}.

\item[{[236]}] RDAL Homepage. \url{https://mem4csd.telecom-paristech.fr/blog/index.php/rdal/}.

\item[{[237]}] W. Reisig. Petri Nets: An Introduction. Springer-Verlag New York, Inc., New York, NY, USA, 1985.

\item[{[238]}] W. Reisig. Petri Nets and Algebraic Specifications. In Theoretical Computer Science, volume 80, pages 1–34. Elsevier, 1991.

\item[{[239]}] W. Reisig. Understanding Petri Nets: Modeling Techniques, Analysis Methods, Case Studies. Springer Berlin Heidelberg, 2013.

\item[{[240]}] J. E. Rivera, D. Ruiz-Gonzalez, F. Lopez-Romero, J. Bautista, and A. Vallecillo. Orchestrating ATL Model Transformations. Proc. of MtATL, 9:34–46, 2009.

\item[{[241]}] C. Rolland. A Comprehensive View of Process Engineering. In International Conference on Advanced Information Systems Engineering, pages 1–24. Springer, 1998.

\item[{[242]}] R. C. Rosenberg. A User’s Guide to ENPORT 4. Wiley, New York, 1974.

\item[{[243]}] J. Rothenberg. The nature of modeling. The Rand Corporation Santa Monica, 1989.

\item[{[244]}] W. W. Royce. Managing the development of large software systems: concepts and techniques. Proc. IEEE WESTCON, Los Angeles, pages 1–9, August 1970. Reprinted in Proceedings of the Ninth International Conference on Software Engineering, March 1987, pp. 328–338.

\item[{[245]}] K. Y. Rozier and M. Y. Vardi. Ltl satisfiability checking. In Model Checking Software, pages 149–167. Springer, 2007.

\item[{[246]}] N. B. Ruparelia. Software Development Lifecycle Models. SIGSOFT Softw. Eng. Notes, 35(3):8–13, May 2010.

\item[{[247]}] SAE AADL Specification. \url{http://standards.sae.org/as5506b/}.

\item[{[248]}] P. Sahlin, A. Bring, and E. Sowell. The Neutral Model Format for Building Simulation, Version 3.02. Royal Institute of Technology, Stockholm, Sweden, 1996.

\item[{[249]}] O. Salo and P. Abrahamsson. Agile methods in european embedded software development organisations: a survey on the actual use and usefulness of extreme programming and scrum. Software, IET, 2(1):58–64, February 2008.

\item[{[250]}] A. Sandholm, P. Fritzson, V. Arora, S. Delp, G. Petersson, and J. Rose. The gait e-book - development of effective participatory learning using simulation and active electronic books. In Proceedings of the 11th Mediterranean Conference on Medical and Biological Engineering and Computing (Medicon’2007), Ljubljana, Slovenia, 2007.

\item[{[251]}] H. Sargent and A. Westerberg. Speed-Up in Chemical Engineering Design. Transaction Institute in Chemical Engineering, 1964.

\item[{[252]}] H. S. Sarjoughian and Y. Chen. Standardizing DEVS models: an endogenous standpoint. In Proceedings of the 2011 Spring Simulation Multiconference, pages 266–273, 2011.

\item[{[253]}] SAVI Homepage. \url{https://savi.avsi.aero/}.

\item[{[254]}] B. Schätz, A. Pretschner, F. Huber, and J. Philipps. Model-based development of embedded systems. In International Conference on Object-Oriented Information Systems, pages 298–311. Springer, 2002.

\item[{[255]}] G. Schmitz. Proceedings of the 4th International Modelica Conference. Technical University Hamburg-Harburg, 2005.

\item[{[256]}] H. Stachowiak. Allgemeine Modelltheorie. 1973.

\item[{[257]}] U. Stern and D. L. Dill. Parallelizing the Mur$\phi$ verifier. In Computer Aided Verification. 9th International Conference, pages 256–267. Springer, 1997.

\item[{[258]}] J. Strauss, D. Augustin, M. Fineberg, B. Johnson, R. Linebarger, and F. Sansom. The sci continuous system simulation language (cssl). SIMULATION, 9(6):281–303, 1967.

\item[{[259]}] E. Syriani, H. Vangheluwe, R. Mannadiar, C. Hansen, S. Van Mierlo, and H. Ergin. AToMPM: A web-based modeling environment. In Joint proceedings of MODELS’13 Invited Talks, Demonstration Session, Poster Session, and ACM Student Research Competition co-located with the 16th International Conference on Model Driven Engineering Languages and Systems (MODELS 2013): September 29-October 4, 2013, Miami, USA, pages 21–25, 2013.

\item[{[260]}] C. Szyperski. Component Software- Beyond Object-Oriented Programming. Addison-Wesley, 1997.

\item[{[261]}] H. Takeuchi and I. Nonaka. The new new product development game. Harvard Business Review, 1986.

\item[{[262]}] R. Tarjan. Depth-first search and linear graph algorithms. In 12th Annual Symposium on Switching and Automata Theory (swat 1971), volume 1, East Lansing, MI, USA, oct 1971.

\item[{[263]}] R. Tarjan. Depth-first search and linear graph algorithms. SIAM journal on computing, 1(2):146–160, 1972.

\item[{[264]}] A. Tarski. A lattice-theoretical fixpoint theorem and its applications. Pacific J. Math., 5(2):285–309, 1955.

\item[{[265]}] E. Technologies. SCADE Suite, 2012. \url{http://www.esterel-technologies.com/products/scade-suite/}.

\item[{[266]}] Y. V. Tendeloo. A Foundation for Multi-Paradigm Modelling. PhD thesis, University of Antwerp, 2018.

\item[{[267]}] Y. Thierry-Mieg. Symbolic model-checking using its-tools. In C. Baier and C. Tinelli, editors, Tools and Algorithms for the Construction and Analysis of Systems - 21st International Conference, TACAS 2015, Held as Part of the European Joint Conferences on Theory and Practice of Software, ETAPS 2015, London, UK, April 11-18, 2015. Proceedings, volume 9035 of Lecture Notes in Computer Science, pages 231–237. Springer, 2015.

\item[{[268]}] J. Thoma. Introduction to Bond Graphs and Their Applications. Pergamon Press, 1975.

\item[{[269]}] J. Thoma. Simulation by Bond Graphs - Introduction to a Graphical Mathod. Springer Verlag, 1989.

\item[{[270]}] M. Tiller. Introduction to Physical Modeling with Modelica. Kluwer Academic Publishers, 2011.

\item[{[271]}] M. Torabzadeh-Tari, Z. Hossain, P. Fritzson, and T. Richter. Omweb - virtual web-based remote laboratory for modelica in engineering courses. In Proceedings of the 8th International Modelica Conference (Modelica’2011), Dresden, Germany, March, 2011.

\item[{[272]}] S. Tripakis. Bridging the semantic gap between heterogeneous modeling formalisms and FMI. In 2015 International Conference on Embedded Computer Systems: Architectures, Modeling, and Simulation (SAMOS), pages 60–69. IEEE, 7 2015.

\item[{[273]}] A. Troccoli and G. Wainer. Implementing Parallel Cell-DEVS. In Proceedings of the 2003 Spring Simulation Symposium, pages 273–280, 2003.

\item[{[274]}] TTool Homepage. \url{https://ttool.telecom-paristech.fr/}.

\item[{[275]}] A. M. Uhrmacher. Dynamic structures in modeling and simulation: a reflective approach. ACM Transactions on Modeling and Computer Simulation, 11:206–232, 2001.

\item[{[276]}] B. Unhelkar. The art of agile practice: A composite approach for projects and organizations. Auerbach Publications, 2016.

\item[{[277]}] A. University. UPPAAL homepage.

\item[{[278]}] A. Valmari. A stubborn attack on state explosion. Form. Methods Syst. Des., 1(4):297–322, 1992.

\item[{[279]}] J. van Amerongen. Dynamical Systems for Creative Technology. Controllab Products B.V., Enschede, 2010.

\item[{[280]}] P. P. J. Van den Bosch and P. Bruijn. The directed digital computer as a teaching tool in control engineering; interactive instruction and design. In Proceedings of the IFAC Symposium on Trends in Automatic Control Education, pages 260–271, 1977.

\item[{[281]}] S. Van Mierlo, Y. Van Tendeloo, B. Meyers, J. Exelmans, and H. Vangheluwe. SCCD: SCXML extended with class diagrams. In 3rd Workshop on Engineering Interactive Systems with SCXML, part of EICS 2016, 2016.

\item[{[282]}] Y. Van Tendeloo and H. Vangheluwe. The modular architecture of the Python(P)DEVS simulation kernel. In Proceedings of the 2014 Spring Simulation Multiconference, pages 387–392, 2014.

\item[{[283]}] Y. Van Tendeloo and H. Vangheluwe. An overview of PythonPDEVS. In C. W. RED, editor, JDF 2016 – Les Journées DEVS Francophones – Théorie et Applications, pages 59–66, 2016.

\item[{[284]}] Y. Van Tendeloo and H. Vangheluwe. An evaluation of DEVS simulation tools. SIMULATION, 93(2):103– 121, 2017.

\item[{[285]}] Y. Van Tendeloo and H. Vangheluwe. Extending the DEVS formalism with initialization information. ArXiv e-prints, 2018.

\item[{[286]}] G. C. V. Vangheluwe, Hans L. and E. J. J. Kerckhoffs. Simulation for the future: Progress of the esprit basic research working group 8467. Proceedings of the 8th European Simulation Symposium, 1996.

\item[{[287]}] H. Vangheluwe. DEVS as a common denominator for multi-formalism hybrid systems modelling. In IEEE International Symposium on Computer-Aided Control System Design, pages 129–134, 2000.

\item[{[288]}] H. Vangheluwe and J. De Lara. Computer automated multi-paradigm modelling for analysis and design of traffic networks. In Proceedings of the 36th conference on Winter simulation, pages 249–258. Winter Simulation Conference, 2004.

\item[{[289]}] H. Vangheluwe, J. De Lara, and P. J. Mosterman. An introduction to multi-paradigm modelling and simulation. In Proceedings of the AIS’2002 conference (AI, Simulation and Planning in High Autonomy Systems), Lisboa, Portugal, pages 9–20, 2002.

\item[{[290]}] H. Vangheluwe, B. H. Li, Y. Reddy, and G. C. Vansteenkiste. A framework for concurrent simulation engineering. Simulation in Concurrent Engineering, pages 50–55, 1993.

\item[{[291]}] H. Vangheluwe, G. Vansteenkiste, and R. Reddy. Simformatics: Meaningful Model Storage, Re-use, Exchange and Simulation. In Summer Computer Simulation Conference, pages 936–944. SOCIETY FOR COMPUTER SIMULATION, ETC, 1997.

\item[{[292]}] H. L. Vangheluwe. DEVS as a common denominator for multi-formalism hybrid systems modelling. In Cacsd. conference proceedings. IEEE international symposium on computer-aided control system design (cat. no. 00th8537), pages 129–134. IEEE, 2000.

\item[{[293]}] H. L. Vangheluwe and G. C. Vansteenkiste. A multi-paradigm modelling and simulation methodology: Formalisms and languages. In Simulation in Industry, Proceedings 8th European Simulation Symposium (Genoa Italy, pages 168–172, 1996.

\item[{[294]}] G. C. Vansteenkiste and H. L. Vangheluwe. European thoughts, actions, and plans for more effective modelling and simulation simulation in europe: A forum for basic research in modelling and simulation. SIMULATION, 66(5):331–335, 1996.

\item[{[295]}] G. C. Vansteenkiste and H. L. Vangheluwe. European Thoughts, Actions, and Plans for More Effective Modelling and Simulation Simulation in Europe: A Forum for Basic Research in Modelling and Simulation. Simulation, 66(5):331–335, 1996.

\item[{[296]}] M. Y. Vardi and P. Wolper. An Automata-Theoretic Approach to Automatic Program Verification (Preliminary Report). In Proceedings 1st Annual IEEE Symp.\textbackslash{} on Logic in Computer Science, LICS’86, Cambridge, MA, USA, 16–18 June 1986, pages 332–344. IEEE, Washington, DC, 1986.

\item[{[297]}] J. Vautherin. Parallel systems specitications with coloured petri nets and algebraic specifications. In G. Rozenberg, editor, Advances in Petri Nets 1987, covers the 7th European Workshop on Applications and Theory of Petri Nets, Oxford, UK, June 1986, volume 266 of Lecture Notes in Computer Science, pages 293–308. Springer, 1986.

\item[{[298]}] Verimag. Openkronos -a model-checker for timed (buchi) automata. \url{http://www-verimag.imag.fr/~tripakis/openkronos.html}.

\item[{[299]}] L. Viklund and P. Fritzson. Objectmath–an object-oriented language and environment for symbolic and numerical processing in scientific computing. Scientific Programming, 4(4):229–250, 1995.

\item[{[300]}] J. Voron, C. Demoulins, and F. Kordon. Adaptable intrusion detection systems dedicated to concurrent programs: A petri net-based approach. In L. Gomes, V. Khomenko, and J. M. Fernandes, editors, 10th International Conference on Application of Concurrency to System Design, ACSD 2010, Braga, Portugal, 21-25 June 2010, pages 57–66. IEEE Computer Society, 2010.

\item[{[301]}] G. A. Wainer. Discrete-Event Modeling and Simulation: A Practitioner’s Approach. CRC Press, 1st edition, 2009.

\item[{[302]}] G. Wanner and E. Hairer. Solving ordinary differential equations I: Nonstiff Problems, volume 1. Springer-Verlag, springer s edition, 1991.

\item[{[303]}] Wikipedia. Capability Maturity Model. \url{https://en.wikipedia.org/wiki/Capability_} Maturity\_Model. Last access: April 2019.

\item[{[304]}] J. C. Willems. The behavioral appraoch to open interconnected systems: Modelling by tearing, zooming, and linking. IEEE Control Systems Magazine, 27:46 – 99, 2007.

\item[{[305]}] S. Wolfram. Statistical mechanics of cellular automata. Reviews of Modern Physics, 55(3):601 – 644, 1983.

\item[{[306]}] S. Wolfram. The Mathematica Book. Wolfram Media Inc., 1997.

\item[{[307]}] T. Zafeer Asici, B. Karaduman, R. Eslampanah, M. Challenger, J. Denil, and H. Vangheluwe. Applying Model Driven Engineering Techniques to the Development of Contiki-based IoT Systems. In 1st International Workshop on Software Engineering Research \& Practices for the Internet of Things (SERP4IoT 2019) held at 41st ACM/IEEE International Conference on Software Engineering (ICSE 2019), Accepted.

\item[{[308]}] B. P. Zeigler, H. Praehofer, and T. G. Kim. Academic Press, 2000.

\item[{[309]}] B. P. Zeigler, H. Praehofer, and T. G. Kim. Theory of Modeling and Simulation. Academic Press, 2nd edition, 2000.

\item[{[310]}] F. Zhang, M. Yeddanapudi, and P. J. Mosterman. Zero-Crossing Location and Detection Algorithms For Hybrid System Simulation. In IFAC Proceedings Volumes, volume 41, pages 7967–7972, Seoul, Korea, jul 2008. Elsevier Ltd.

\end{description}
